\chapter{Functions in Python: Reusable Blocks of Code}

\section{Introduction: Real-Life Hook}

Students use routines in daily life:  

\begin{itemize}
    \item Brushing teeth every morning  
    \item Sending birthday wishes  
    \item Calculating total marks  
\end{itemize}

Instead of repeating the same steps, we create \textbf{functions} in Python to \textbf{reuse code easily}.

\section{Concept: What is a Function?}

\begin{definition}{Function}
A \textbf{function} is a block of code that performs a specific task. It can be reused multiple times in a program.
\end{definition}

\subsection{Defining a Function}

\begin{verbatim}
def greet():
    print("Hello! Welcome to Python class")
\end{verbatim}

- \texttt{def} keyword defines the function  
- \texttt{greet()} is the function name  
- Indented block is the function body

\subsection{Calling a Function}

\begin{verbatim}
greet()
greet()
\end{verbatim}

Output:

\begin{verbatim}
Hello! Welcome to Python class
Hello! Welcome to Python class
\end{verbatim}

\section{Functions with Parameters}

Functions can take input values (parameters):

\begin{verbatim}
def greet(name):
    print("Hello,", name)

greet("Amina")
greet("Rafi")
\end{verbatim}

Output:

\begin{verbatim}
Hello, Amina
Hello, Rafi
\end{verbatim}

\section{Functions with Return Values}

Functions can return values for later use:

\begin{verbatim}
def add(a, b):
    return a + b

result = add(5, 3)
print("Sum is:", result)
\end{verbatim}

Output:

\begin{verbatim}
Sum is: 8
\end{verbatim}

\section{Step-by-Step Diagram: Function Flow}

\begin{center}
\begin{tikzpicture}[node distance=1.5cm]
\node (start) [draw, ellipse, fill=green!20] {Start Program};
\node (define) [draw, rectangle, fill=yellow!20, below of=start] {Define Function};
\node (call) [draw, rectangle, fill=orange!20, below of=define] {Call Function};
\node (execute) [draw, rectangle, fill=orange!20, below of=call] {Execute Function Body};
\node (return) [draw, rectangle, fill=yellow!20, below of=execute] {Return Value (Optional)};
\node (end) [draw, ellipse, fill=green!20, below of=return] {End Program};
\draw[->] (start) -- (define);
\draw[->] (define) -- (call);
\draw[->] (call) -- (execute);
\draw[->] (execute) -- (return);
\draw[->] (return) -- (end);
\end{tikzpicture}
\end{center}

\section{Mini-Project: Calculator with Functions}

Create a program with functions for addition, subtraction, multiplication, and division:

\begin{verbatim}
def add(a, b):
    return a + b

def subtract(a, b):
    return a - b

def multiply(a, b):
    return a * b

def divide(a, b):
    return a / b

num1 = int(input("Enter first number: "))
num2 = int(input("Enter second number: "))

print("Addition:", add(num1, num2))
print("Subtraction:", subtract(num1, num2))
print("Multiplication:", multiply(num1, num2))
print("Division:", divide(num1, num2))
\end{verbatim}

\section{Using Functions to Organize Code}

- Functions help avoid repeating code  
- Functions make programs readable and maintainable  
- Complex programs are easier to debug when using functions

\section{Key Concepts}

\begin{itemize}
    \item A function is a reusable block of code  
    \item Functions can take parameters (inputs)  
    \item Functions can return values  
    \item Functions improve code readability, reduce repetition, and organize programs  
\end{itemize}

\section{End-of-Chapter Exercises}

\begin{exercise}{Very Short Questions}
\begin{enumerate}
    \item What is a function in Python?  
    \item Which keyword is used to define a function?  
    \item Can a function return a value?  
\end{enumerate}
\end{exercise}

\begin{exercise}{Short Answer Questions}
\begin{enumerate}
    \item Write a function to print your name.  
    \item Write a function that takes two numbers and returns their product.  
    \item Explain the difference between parameters and return values.  
\end{enumerate}
\end{exercise}

\begin{exercise}{Activity-Based Questions}
\begin{enumerate}
    \item Create a program with functions for addition, subtraction, multiplication, and division of two numbers.  
    \item Write a function that takes a list of numbers and returns the average.  
    \item Modify the calculator program to use a menu inside a function to choose the operation.  
\end{enumerate}
\end{exercise}

\section{Teacher Notes}

\marginnote{
\textbf{Learning Objectives:}  
\begin{itemize}
    \item Students understand how to define and call functions.  
    \item Students can pass parameters and return values.  
    \item Students can organize programs into reusable blocks.  
\end{itemize}

\textbf{Teaching Tips:}  
\begin{itemize}
    \item Demonstrate simple functions first, then parameters and return values.  
    \item Show real-life analogy: “function as a routine you do repeatedly”.  
\end{itemize}

\textbf{Assessment:}  
\begin{itemize}
    \item Check if students can create and call functions correctly.  
    \item Observe understanding of parameters, return values, and code organization.  
\end{itemize}
}